\IfFileExists{stacks-project.cls}{%
\documentclass{stacks-project}
}{%
\documentclass{amsart}
}

% For dealing with references we use the comment environment
\usepackage{verbatim}
\newenvironment{reference}{\comment}{\endcomment}
%\newenvironment{reference}{}{}
\newenvironment{slogan}{\comment}{\endcomment}
\newenvironment{history}{\comment}{\endcomment}

% For commutative diagrams we use Xy-pic
\usepackage[all]{xy}

% We use 2cell for 2-commutative diagrams.
\xyoption{2cell}
\UseAllTwocells

% We use multicol for the list of chapters between chapters
\usepackage{multicol}

% This is generally recommended for better output
\usepackage{lmodern}
\usepackage[T1]{fontenc}

% For cross-file-references
\usepackage{xr-hyper}

% Package for hypertext links:
\usepackage{hyperref}

% For any local file, say "hello.tex" you want to link to please
% use \externaldocument[hello-]{hello}
\externaldocument[introduction-]{introduction}
\externaldocument[conventions-]{conventions}
\externaldocument[sets-]{sets}
\externaldocument[categories-]{categories}
\externaldocument[topology-]{topology}
\externaldocument[sheaves-]{sheaves}
\externaldocument[sites-]{sites}
\externaldocument[stacks-]{stacks}
\externaldocument[fields-]{fields}
\externaldocument[algebra-]{algebra}
\externaldocument[brauer-]{brauer}
\externaldocument[homology-]{homology}
\externaldocument[derived-]{derived}
\externaldocument[simplicial-]{simplicial}
\externaldocument[more-algebra-]{more-algebra}
\externaldocument[smoothing-]{smoothing}
\externaldocument[modules-]{modules}
\externaldocument[sites-modules-]{sites-modules}
\externaldocument[injectives-]{injectives}
\externaldocument[cohomology-]{cohomology}
\externaldocument[sites-cohomology-]{sites-cohomology}
\externaldocument[dga-]{dga}
\externaldocument[dpa-]{dpa}
\externaldocument[sdga-]{sdga}
\externaldocument[hypercovering-]{hypercovering}
\externaldocument[schemes-]{schemes}
\externaldocument[constructions-]{constructions}
\externaldocument[properties-]{properties}
\externaldocument[morphisms-]{morphisms}
\externaldocument[coherent-]{coherent}
\externaldocument[divisors-]{divisors}
\externaldocument[limits-]{limits}
\externaldocument[varieties-]{varieties}
\externaldocument[topologies-]{topologies}
\externaldocument[descent-]{descent}
\externaldocument[perfect-]{perfect}
\externaldocument[more-morphisms-]{more-morphisms}
\externaldocument[flat-]{flat}
\externaldocument[groupoids-]{groupoids}
\externaldocument[more-groupoids-]{more-groupoids}
\externaldocument[etale-]{etale}
\externaldocument[chow-]{chow}
\externaldocument[intersection-]{intersection}
\externaldocument[pic-]{pic}
\externaldocument[weil-]{weil}
\externaldocument[adequate-]{adequate}
\externaldocument[dualizing-]{dualizing}
\externaldocument[duality-]{duality}
\externaldocument[discriminant-]{discriminant}
\externaldocument[derham-]{derham}
\externaldocument[local-cohomology-]{local-cohomology}
\externaldocument[algebraization-]{algebraization}
\externaldocument[curves-]{curves}
\externaldocument[resolve-]{resolve}
\externaldocument[models-]{models}
\externaldocument[functors-]{functors}
\externaldocument[equiv-]{equiv}
\externaldocument[pione-]{pione}
\externaldocument[etale-cohomology-]{etale-cohomology}
\externaldocument[proetale-]{proetale}
\externaldocument[relative-cycles-]{relative-cycles}
\externaldocument[more-etale-]{more-etale}
\externaldocument[trace-]{trace}
\externaldocument[crystalline-]{crystalline}
\externaldocument[spaces-]{spaces}
\externaldocument[spaces-properties-]{spaces-properties}
\externaldocument[spaces-morphisms-]{spaces-morphisms}
\externaldocument[decent-spaces-]{decent-spaces}
\externaldocument[spaces-cohomology-]{spaces-cohomology}
\externaldocument[spaces-limits-]{spaces-limits}
\externaldocument[spaces-divisors-]{spaces-divisors}
\externaldocument[spaces-over-fields-]{spaces-over-fields}
\externaldocument[spaces-topologies-]{spaces-topologies}
\externaldocument[spaces-descent-]{spaces-descent}
\externaldocument[spaces-perfect-]{spaces-perfect}
\externaldocument[spaces-more-morphisms-]{spaces-more-morphisms}
\externaldocument[spaces-flat-]{spaces-flat}
\externaldocument[spaces-groupoids-]{spaces-groupoids}
\externaldocument[spaces-more-groupoids-]{spaces-more-groupoids}
\externaldocument[bootstrap-]{bootstrap}
\externaldocument[spaces-pushouts-]{spaces-pushouts}
\externaldocument[spaces-chow-]{spaces-chow}
\externaldocument[groupoids-quotients-]{groupoids-quotients}
\externaldocument[spaces-more-cohomology-]{spaces-more-cohomology}
\externaldocument[spaces-simplicial-]{spaces-simplicial}
\externaldocument[spaces-duality-]{spaces-duality}
\externaldocument[formal-spaces-]{formal-spaces}
\externaldocument[restricted-]{restricted}
\externaldocument[spaces-resolve-]{spaces-resolve}
\externaldocument[formal-defos-]{formal-defos}
\externaldocument[defos-]{defos}
\externaldocument[cotangent-]{cotangent}
\externaldocument[examples-defos-]{examples-defos}
\externaldocument[algebraic-]{algebraic}
\externaldocument[examples-stacks-]{examples-stacks}
\externaldocument[stacks-sheaves-]{stacks-sheaves}
\externaldocument[criteria-]{criteria}
\externaldocument[artin-]{artin}
\externaldocument[quot-]{quot}
\externaldocument[stacks-properties-]{stacks-properties}
\externaldocument[stacks-morphisms-]{stacks-morphisms}
\externaldocument[stacks-limits-]{stacks-limits}
\externaldocument[stacks-cohomology-]{stacks-cohomology}
\externaldocument[stacks-perfect-]{stacks-perfect}
\externaldocument[stacks-introduction-]{stacks-introduction}
\externaldocument[stacks-more-morphisms-]{stacks-more-morphisms}
\externaldocument[stacks-geometry-]{stacks-geometry}
\externaldocument[moduli-]{moduli}
\externaldocument[moduli-curves-]{moduli-curves}
\externaldocument[examples-]{examples}
\externaldocument[exercises-]{exercises}
\externaldocument[guide-]{guide}
\externaldocument[desirables-]{desirables}
\externaldocument[coding-]{coding}
\externaldocument[obsolete-]{obsolete}
\externaldocument[fdl-]{fdl}
\externaldocument[index-]{index}

% Theorem environments.
%
\theoremstyle{plain}
\newtheorem{theorem}[subsection]{Theorem}
\newtheorem{proposition}[subsection]{Proposition}
\newtheorem{lemma}[subsection]{Lemma}

\theoremstyle{definition}
\newtheorem{definition}[subsection]{Definition}
\newtheorem{example}[subsection]{Example}
\newtheorem{exercise}[subsection]{Exercise}
\newtheorem{situation}[subsection]{Situation}

\theoremstyle{remark}
\newtheorem{remark}[subsection]{Remark}
\newtheorem{remarks}[subsection]{Remarks}

\numberwithin{equation}{subsection}

% Macros
%
\def\lim{\mathop{\mathrm{lim}}\nolimits}
\def\colim{\mathop{\mathrm{colim}}\nolimits}
\def\Spec{\mathop{\mathrm{Spec}}}
\def\Hom{\mathop{\mathrm{Hom}}\nolimits}
\def\Ext{\mathop{\mathrm{Ext}}\nolimits}
\def\SheafHom{\mathop{\mathcal{H}\!\mathit{om}}\nolimits}
\def\SheafExt{\mathop{\mathcal{E}\!\mathit{xt}}\nolimits}
\def\Sch{\mathit{Sch}}
\def\Mor{\mathop{\mathrm{Mor}}\nolimits}
\def\Ob{\mathop{\mathrm{Ob}}\nolimits}
\def\Sh{\mathop{\mathit{Sh}}\nolimits}
\def\NL{\mathop{N\!L}\nolimits}
\def\CH{\mathop{\mathrm{CH}}\nolimits}
\def\proetale{{pro\text{-}\acute{e}tale}}
\def\etale{{\acute{e}tale}}
\def\QCoh{\mathit{QCoh}}
\def\Ker{\mathop{\mathrm{Ker}}}
\def\Im{\mathop{\mathrm{Im}}}
\def\Coker{\mathop{\mathrm{Coker}}}
\def\Coim{\mathop{\mathrm{Coim}}}

% Boxtimes
%
\DeclareMathSymbol{\boxtimes}{\mathbin}{AMSa}{"02}

%
% Macros for moduli stacks/spaces
%
\def\QCohstack{\mathcal{QC}\!\mathit{oh}}
\def\Cohstack{\mathcal{C}\!\mathit{oh}}
\def\Spacesstack{\mathcal{S}\!\mathit{paces}}
\def\Quotfunctor{\mathrm{Quot}}
\def\Hilbfunctor{\mathrm{Hilb}}
\def\Curvesstack{\mathcal{C}\!\mathit{urves}}
\def\Polarizedstack{\mathcal{P}\!\mathit{olarized}}
\def\Complexesstack{\mathcal{C}\!\mathit{omplexes}}
% \Pic is the operator that assigns to X its picard group, usage \Pic(X)
% \Picardstack_{X/B} denotes the Picard stack of X over B
% \Picardfunctor_{X/B} denotes the Picard functor of X over B
\def\Pic{\mathop{\mathrm{Pic}}\nolimits}
\def\Picardstack{\mathcal{P}\!\mathit{ic}}
\def\Picardfunctor{\mathrm{Pic}}
\def\Deformationcategory{\mathcal{D}\!\mathit{ef}}
\usepackage{cleveref}
\crefname{originalproposition}{original proposition}{original propositions}
\Crefname{originalproposition}{Original proposition}{Original propositions}
\makeatletter
\newlabel{limits-proposition-approximate}{{01ZA}{source}{}{sourcepointer-1}{}}
\newlabel{limits-proposition-approximate@cref}{{[originalproposition][0][]01ZA}{[1][0][]source}{}{}{}}
\makeatother
\crefname{originallemma}{original lemma}{original lemmas}
\Crefname{originallemma}{Original lemma}{Original lemmas}
\makeatletter
\newlabel{lemma-resolution-property-descends}{{0F8B}{source}{}{sourcepointer-2}{}}
\newlabel{lemma-resolution-property-descends@cref}{{[originallemma][0][]0F8B}{[1][0][]source}{}{}{}}
\makeatother
\crefname{originalsection}{original section}{original sections}
\Crefname{originalsection}{Original section}{Original sections}
\makeatletter
\newlabel{schemes-section-fibre-products}{{01JO}{source}{}{sourcepointer-3}{}}
\newlabel{schemes-section-fibre-products@cref}{{[originalsection][0][]01JO}{[1][0][]source}{}{}{}}
\makeatother
\crefname{originallemma}{original lemma}{original lemmas}
\Crefname{originallemma}{Original lemma}{Original lemmas}
\makeatletter
\newlabel{morphisms-lemma-characterize-affine}{{01S8}{source}{}{sourcepointer-4}{}}
\newlabel{morphisms-lemma-characterize-affine@cref}{{[originallemma][0][]01S8}{[1][0][]source}{}{}{}}
\makeatother
\crefname{originallemma}{original lemma}{original lemmas}
\Crefname{originallemma}{Original lemma}{Original lemmas}
\makeatletter
\newlabel{coherent-lemma-criterion-affine-morphism}{{0F83}{source}{}{sourcepointer-5}{}}
\newlabel{coherent-lemma-criterion-affine-morphism@cref}{{[originallemma][0][]0F83}{[1][0][]source}{}{}{}}
\makeatother
\crefname{originalproposition}{original proposition}{original propositions}
\Crefname{originalproposition}{Original proposition}{Original propositions}
\makeatletter
\newlabel{proposition-Noetherian}{{09T4}{source}{}{sourcepointer-6}{}}
\newlabel{proposition-Noetherian@cref}{{[originalproposition][0][]09T4}{[1][0][]source}{}{}{}}
\makeatother
\crefname{originallemma}{original lemma}{original lemmas}
\Crefname{originallemma}{Original lemma}{Original lemmas}
\makeatletter
\newlabel{properties-lemma-quasi-coherent-colimit-finite-type}{{01PG}{source}{}{sourcepointer-7}{}}
\newlabel{properties-lemma-quasi-coherent-colimit-finite-type@cref}{{[originallemma][0][]01PG}{[1][0][]source}{}{}{}}
\makeatother
\crefname{originallemma}{original lemma}{original lemmas}
\Crefname{originallemma}{Original lemma}{Original lemmas}
\makeatletter
\newlabel{coherent-lemma-quasi-coherent-affine-cohomology-zero}{{01XB}{source}{}{sourcepointer-8}{}}
\newlabel{coherent-lemma-quasi-coherent-affine-cohomology-zero@cref}{{[originallemma][0][]01XB}{[1][0][]source}{}{}{}}
\makeatother
\crefname{originallemma}{original lemma}{original lemmas}
\Crefname{originallemma}{Original lemma}{Original lemmas}
\makeatletter
\newlabel{cohomology-lemma-dual-perfect-complex}{{08DQ}{source}{}{sourcepointer-9}{}}
\newlabel{cohomology-lemma-dual-perfect-complex@cref}{{[originallemma][0][]08DQ}{[1][0][]source}{}{}{}}
\makeatother
\title{A quasi-compact quasi-separated scheme with the resolution property has affine diagonal}
\author{The Stacks Project Authors}
\date{stacks-project:perfect.tex:local-snapshot-20261012; GNU FDL 1.2 or later}
\begin{document}
\maketitle
\setcounter{subsection}{1}
\noindent\textit{Editorial note.} All delivered mathematical text is the author's own, copied from the pinned Stacks Project chapter file perfect.tex; no statement or proof step was rewritten, reordered or removed. The counted span is the author environment and its own proof environment. References inside the span are handed to the assembler as navigation stubs that carry the canonical tag of the cited tag only; whether each of them is load-bearing is recorded on the stub itself.\par\medskip
\noindent\textit{Allowed standard prerequisites (editorial).}\begin{itemize}
\item The resolution property descends along approximations, so the statement reduces to Noetherian schemes. invoked at proof lines 9713-9716. Fixes the Noetherian hypothesis used later via the derived category equivalence.
\item A morphism is affine as soon as the inverse images of the affine opens of the target are affine, equivalently R\textasciicircum{}1 of the pushforward of every quasi-coherent module vanishes. invoked at proof lines 9722-9730. Reduces affinity of the diagonal to the vanishing of R\textasciicircum{}1j\_* for inclusions of affine opens.
\item On a Noetherian scheme the derived pushforward of a quasi-coherent module is represented by a complex of quasi-coherent modules, and every quasi-coherent module is the filtered union of its coherent submodules. invoked at proof lines 9733-9748. Lets a hypothetical nonzero H\textasciicircum{}1 be detected by a coherent module F mapping nontrivially to it.
\item A nonzero map to the cohomology of Rj\_*G corresponds by adjunction and duality with a finite locally free module to a nonzero class in the Cech cohomology H\textasciicircum{}1 of an open, which vanishes over affines. invoked at proof lines 9750-9775. Produces the contradiction from the resolution property.
\end{itemize}\medskip
\ifdefined\hypertarget\hypertarget{sourcepointer-1}{}\fi
\noindent\textit{Source reference.} Original proposition 01ZA: sources/stacks-chow/tags.txt, source lines 1--1. Cited at proof lines 9713-9715 for the reduction to the Noetherian case; the descent step and the whole Noetherian argument are inside the counted span.\par\medskip
\ifdefined\hypertarget\hypertarget{sourcepointer-2}{}\fi
\noindent\textit{Source reference.} Original lemma 0F8B: sources/stacks-perfect/pinned-original-source.tex, source lines 9661--9661. Cited at proof lines 9714-9716 together with the approximation proposition; the Noetherian reduction is completed inside the counted span.\par\medskip
\ifdefined\hypertarget\hypertarget{sourcepointer-3}{}\fi
\noindent\textit{Source reference.} Original section 01JO: sources/stacks-chow/tags.txt, source lines 1--1. Cited at proof lines 9719-9721 for the affine open cover of X x X; the reduction of affinity of the diagonal is inside the counted span.\par\medskip
\ifdefined\hypertarget\hypertarget{sourcepointer-4}{}\fi
\noindent\textit{Source reference.} Original lemma 01S8: sources/stacks-chow/tags.txt, source lines 1--1. Cited at proof lines 9722-9724 to reduce to U  \textbackslash{}cap  V affine; the further reduction to R\textasciicircum{}1j\_* = 0 is inside the counted span.\par\medskip
\ifdefined\hypertarget\hypertarget{sourcepointer-5}{}\fi
\noindent\textit{Source reference.} Original lemma 0F83: sources/stacks-chow/tags.txt, source lines 1--1. Cited at proof lines 9726-9730 to reduce to R\textasciicircum{}1j\_*G = 0; the contradiction argument against nonzero R\textasciicircum{}1 is inside the counted span.\par\medskip
\ifdefined\hypertarget\hypertarget{sourcepointer-6}{}\fi
\noindent\textit{Source reference.} Original proposition 09T4: sources/stacks-perfect/pinned-original-source.tex, source lines 2026--2026. Cited at proof lines 9733-9736 to represent Rj\_*G by H\textasciicircum{}bullet; the contradiction construction on the cohomology of H\textasciicircum{}bullet is inside the counted span.\par\medskip
\ifdefined\hypertarget\hypertarget{sourcepointer-7}{}\fi
\noindent\textit{Source reference.} Original lemma 01PG: sources/stacks-chow/tags.txt, source lines 1--1. Cited at proof lines 9744-9748 to pick the coherent F detecting H\textasciicircum{}1; the choice of the surjection from E is inside the counted span.\par\medskip
\ifdefined\hypertarget\hypertarget{sourcepointer-8}{}\fi
\noindent\textit{Source reference.} Original lemma 01XB: sources/stacks-chow/tags.txt, source lines 1--1. Cited at proof lines 9770-9773 to produce the contradiction; the adjunction and duality steps leading to the H\textasciicircum{}1 class are inside the counted span.\par\medskip
\ifdefined\hypertarget\hypertarget{sourcepointer-9}{}\fi
\noindent\textit{Source reference.} Original lemma 08DQ: sources/stacks-chow/tags.txt, source lines 1--1. Cited at proof lines 9773-9775 as a reference for the dualization step; the dualization itself and the conclusion are inside the counted span.\par\medskip
\begin{lemma}
\label{lemma-resolution-property-affine-diagonal}
\begin{reference}
Special case of \cite[Proposition 1.3]{totaro_resolution}.
\end{reference}
Let $X$ be a quasi-compact and quasi-separated scheme with
the resolution property. Then $X$ has affine diagonal.
\end{lemma}

\begin{proof}
Combining Limits, Proposition \ref{limits-proposition-approximate}
and Lemma \ref{lemma-resolution-property-descends}
this reduces to the case where $X$ is Noetherian (small detail omitted).
Assume $X$ is Noetherian.
Recall that $X \times X$ is covered by the affine opens
$U \times V$ for affine opens $U$, $V$ of $X$, see
Schemes, Section \ref{schemes-section-fibre-products}.
Hence to show that the diagonal $\Delta : X \to X \times X$
is affine, it suffices to show that $U \cap V = \Delta^{-1}(U \times V)$
is affine for all affine opens $U$, $V$ of $X$, see
Morphisms, Lemma \ref{morphisms-lemma-characterize-affine}.
In particular, it suffices to show that the inclusion morphism
$j : U \to X$ is affine if $U$ is an affine open of $X$.
By Cohomology of Schemes, Lemma \ref{coherent-lemma-criterion-affine-morphism}
it suffices to show that $R^1j_*\mathcal{G} = 0$ for any
quasi-coherent $\mathcal{O}_U$-module $\mathcal{G}$.
By Proposition \ref{proposition-Noetherian} (this is where we use
that we've reduced to the Noetherian case)
we can represent $Rj_*\mathcal{G}$ by a complex
$\mathcal{H}^\bullet$ of quasi-coherent $\mathcal{O}_X$-modules.
Assume
$$
H^1(\mathcal{H}^\bullet) =
\Ker(\mathcal{H}^1 \to \mathcal{H}^2)/\Im(\mathcal{H}^0 \to \mathcal{H}^1)
$$
is nonzero in order to get a contradiction. Then we can find a coherent
$\mathcal{O}_X$-module $\mathcal{F}$ and a map
$$
\mathcal{F} \longrightarrow
\Ker(\mathcal{H}^1 \to \mathcal{H}^2)
$$
such that the composition with the projection onto $H^1(\mathcal{H}^\bullet)$
is nonzero. Namely, we can write $\Ker(\mathcal{H}^1 \to \mathcal{H}^2)$
as the filtered union of its coherent submodules by
Properties, Lemma \ref{properties-lemma-quasi-coherent-colimit-finite-type}
and then one of these will do the job.
Next, we choose a finite locally free $\mathcal{O}_X$-module
$\mathcal{E}$ and a surjection $\mathcal{E} \to \mathcal{F}$ using
the resolution property of $X$.
This produces a map in the derived category
$$
\mathcal{E}[-1] \longrightarrow Rj_*\mathcal{G}
$$
which is nonzero on cohomology sheaves and hence nonzero in $D(\mathcal{O}_X)$.
By adjunction, this is the same thing as a map
$$
j^*\mathcal{E}[-1] \to \mathcal{G}
$$
nonzero in $D(\mathcal{O}_U)$. Since $\mathcal{E}$ is finite locally
free this is the same thing as a nonzero element of
$$
H^1(U, j^*\mathcal{E}^\vee \otimes_{\mathcal{O}_U} \mathcal{G})
$$
where
$\mathcal{E}^\vee = \SheafHom_{\mathcal{O}_X}(\mathcal{E}, \mathcal{O}_X)$
is the dual finite locally free module. However, this group is zero by
Cohomology of Schemes, Lemma
\ref{coherent-lemma-quasi-coherent-affine-cohomology-zero}
which is the desired contradiction.
(If in doubt about the step using duals, please see the more general
Cohomology, Lemma \ref{cohomology-lemma-dual-perfect-complex}.)
\end{proof}

\end{document}
